% ID: SEL-D0067-OA-E01
% ESTADO: cotejada. Segunda pasada visual y borrador compilado revisados.
% Fuente: Selectividad/sel_2009_matematicas_aplicadas/mat apli CC SS01.pdf
% SHA-256: f5e34036805266809f1b5a2c6087d7ae598afd929b5179ec5ea3888db6cbccf4
% Página: 1; opción A; ejercicio original 1.

\begin{enumerate}[label=\alph*)]
\item \textbf{(1.5 puntos)} En un comercio de bricolaje se venden listones de madera de tres longitudes: $0.90$ m, $1.50$ m y $2.40$ m, cuyos precios respectivos son 4 euros, 6 euros y 10 euros. Un cliente ha comprado 19 listones, con una longitud total de 30 m, que le han costado 126 euros en total.

Plantee, sin resolver, el sistema de ecuaciones necesario para determinar cuántos listones de cada longitud ha comprado este cliente.
\item \textbf{(1.5 puntos)} Clasifique el siguiente sistema de ecuaciones y resuélvalo, si es posible:
\[
\left.\begin{aligned}3x-y-z&=0\\2x-2y+z&=18\\x-3z&=0\end{aligned}\right\}.
\]
\end{enumerate}
